\chapter{More than one objective}
\label{sec:moo}

Several of our workloads do not have a single number to minimize, and the
Hamiltonian network introduced in Section~\ref{sec:running} is the cleanest
case: a model that nails one-step prediction but leaks energy over long
rollouts is not better than one that does the reverse; it is better at one
thing and worse at another. When the exchange rate between objectives is
known in advance, the problem collapses back to
Section~\ref{sec:formal}: fix the weights and minimize the combination.
Usually the exchange rate is exactly what we do not know, and this chapter
builds the machinery for optimizing honestly without inventing one. The
build order mirrors the single-objective story: first the right notion of
``better'' (dominance), then the right notion of progress (hypervolume),
then the acquisition functions and schedulers that optimize it.

\section{Dominance, the Pareto front, and how to score progress}

With objectives $f_1, \dots, f_m$ to minimize, configuration $\config$
\emph{dominates} $\config'$ if it is no worse in every objective and
strictly better in at least one: there is no argument for keeping
$\config'$ when $\config$ is on the table. (Economists are on home
ground in this whole chapter: this is Pareto dominance exactly as
welfare theory uses it, and the scalarization results below are the
optimization mirror of familiar aggregation problems.) Domination is only a partial
order, and that is the entire difficulty of this chapter: most pairs of
configurations are simply incomparable, one better at prediction, the
other at energy conservation, and no amount of cleverness removes the
incomparability without importing a preference.

The \emph{Pareto front} is
the image of the undominated configurations, the trade-offs that survive
every argument (Figure~\ref{fig:pareto}). Surveys of multi-objective
tuning sort methods by when preferences enter: before the search, which
fixes a scalarization and returns one point, or after it, which
approximates the front and lets the user choose with the trade-off in view
\citep{karl2023moo, morales2022survey}. Everything below is of the second
kind unless said otherwise: the tuner returns an approximate front, and
the exchange rate is made visible rather than assumed.

If the target is a set rather than a point, what does progress mean? The
standard scalar score is the \emph{hypervolume}: fix a reference point
$\mathbf{r} \in \R^m$ that is worse than anything acceptable, and measure
the volume of objective space that the current front dominates,
\begin{equation}
\mathrm{HV}(P) \;=\; \operatorname{vol}\!\Big(
\bigcup_{\mathbf{y} \in P} \; [\mathbf{y}, \mathbf{r}] \Big),
\label{eq:hv}
\end{equation}
the union of the axis-aligned boxes spanned between each front point
$\mathbf{y}$ and the reference. The indicator goes back to
\citet{zitzler1999spea}.

A small computation makes it concrete. Take
the Hamiltonian study with prediction error and energy drift both scaled
to $[0,5]$, reference point $\mathbf{r} = (5,5)$, and three undominated
trials at $(1,4)$, $(2,2)$, and $(4,1)$. Slicing the dominated region
into vertical strips: for first-objective values in $[1,2)$ only the
point $(1,4)$ dominates, contributing a strip of height $1$ and area
$1$; on $[2,4)$ the point $(2,2)$ takes over, height $3$, area $6$; on
$[4,5)$ the point $(4,1)$ gives height $4$, area $4$. Total hypervolume:
$11$ of the $25$ available. Adding a new trial at $(3, 1.5)$ would carve
out fresh territory (a rectangle of width $1$ and height $0.5$ that no
existing box covers), so the hypervolume rises; adding a dominated trial
adds nothing. This is why the indicator earns its role as the standard
yardstick: by \eqref{eq:hv} it rewards convergence toward the true front
and spread along it at once, and any addition that dominates new
territory strictly increases it.

\begin{figure}[t]
\centering
\begin{tikzpicture}[scale=1.05, font=\small]
\draw[-{Latex}] (0,0) -- (5.9,0) node[below left=1mm and 0mm] {prediction error $f_1$};
\draw[-{Latex}] (0,0) -- (0,5.9);
\node[rotate=90, anchor=south] at (-0.4,2.9) {energy drift $f_2$};
% reference point
\fill (5,5) circle (2pt) node[above right] {$\mathbf{r}=(5,5)$};
% dominated region (union of boxes), strips
\fill[blue!12] (1,4) rectangle (2,5);
\fill[blue!12] (2,2) rectangle (4,5);
\fill[blue!12] (4,1) rectangle (5,5);
% front points
\fill[blue!70!black] (1,4) circle (2.2pt) node[below left] {$(1,4)$};
\fill[blue!70!black] (2,2) circle (2.2pt) node[below left] {$(2,2)$};
\fill[blue!70!black] (4,1) circle (2.2pt) node[below left] {$(4,1)$};
% a dominated point
\fill[red!70!black] (3.4,3.4) circle (2.2pt);
\node[red!70!black, anchor=west] at (3.5,3.55) {dominated by $(2,2)$};
% front staircase
\draw[thick, blue!70!black] (1,5) -- (1,4) -- (2,4) -- (2,2) -- (4,2) -- (4,1) -- (5,1);
\node[anchor=west, align=left] at (1.35,5.55)
  {shaded area $=$ hypervolume $= 11$};
\end{tikzpicture}
\caption{The Hamiltonian study's objective plane (schematic, minimizing
both axes). The three blue trials are mutually incomparable and form the
current front; the red trial is dominated and argues for nothing. The
shaded region is everything the front dominates up to the reference
point, and its area, computed strip by strip in the text, is the
hypervolume~\eqref{eq:hv}: the one number that grows only when the front
genuinely advances or spreads.}
\label{fig:pareto}
\end{figure}

Hypervolume's cost is the catch. Computing and decomposing the dominated
region has time complexity that grows super-polynomially in the number of
objectives $m$ \citep{daulton2021qnehvi}, which in practice confines
hypervolume-driven methods to two to four objectives; beyond that,
scalarization-based methods take over. All of our current use cases sit
comfortably at $m \in \{2, 3\}$.

\section{Scalarize many times, or evolve a population}

The oldest trick converts the problem into many single-objective ones:
draw a weight vector, optimize the scalarized objective, repeat with
fresh weights, and collect the results. Whether this can even recover the
whole front depends on the scalarization, and the failure of the obvious
choice is geometric enough to draw.

Minimizing a weighted sum
$\sum_i w_i f_i$ with fixed positive weights means sliding the line (in
general, hyperplane) $\sum_i w_i f_i = c$ toward the origin until it last
touches the attainable region: the touching point is the minimizer. A
supporting line can only touch the region's convex hull, so for
\emph{any} positive weights, the minimizer of a weighted sum sits on the
convex hull of the attainable objective vectors, and Pareto points inside
a concave stretch of the front are never selected, whatever the weights
(Figure~\ref{fig:chebyshev}, left).

In the Hamiltonian study this is not
a curiosity: if the best balanced models happen to sit in a dent of the
front, every weighted-sum study we run will report only the two extremes,
a model that predicts well and leaks, and a model that conserves and
mispredicts, with the interesting compromises silently unreachable.

The scalarization of choice is therefore the Chebyshev form: with an
ideal point $\mathbf{z}^{\star}$ (each objective's best value, in
practice the best seen) and weights $w_i > 0$,
\begin{equation}
g_{\mathbf{w}}(\config)
= \max_{i = 1, \dots, m}\; w_i \left( f_i(\config) - z_i^{\star} \right),
\label{eq:chebyshev}
\end{equation}
the weighted \emph{worst} gap to the ideal rather than the weighted sum
of gaps. Its level sets are axis-aligned corners rather than lines, and a
corner can nestle into a concave dent that no line can touch
(Figure~\ref{fig:chebyshev}, right). The standard completeness statement
(sweeping positive weights with an appropriate true ideal point recovers
every Pareto point under convexity or boundary assumptions) is proven
using the true utopian vector, not a best-seen point; the latter is
usable in practice and provides a coverage advantage over weighted sums,
but it is an operational approximation rather than a theoretical
guarantee \citep{karl2023moo}.

ParEGO runs exactly this inside BO: each iteration
draws a random Chebyshev weight, scalarizes the observed data, and takes
an EI step on the result \citep{knowles2006parego};
random-scalarization bandits generalize the idea with guarantees
\citep{paria2019scalarization}. Scalarization scales gracefully in $m$
and composes with every single-objective engine we build, which is why
it remains the fallback when hypervolume machinery gets expensive.

\begin{figure}[t]
\centering
\begin{tikzpicture}[scale=0.82, font=\small]
% ---- left panel: weighted sum ----
\begin{scope}
\draw[-{Latex}] (0,0) -- (5.4,0) node[below left=1mm and 0mm] {$f_1$};
\draw[-{Latex}] (0,0) -- (0,5.4) node[left] {$f_2$};
% attainable region boundary: front with concave dent
\draw[thick, blue!70!black]
  (0.6,4.8) .. controls (0.7,3.4) and (1.0,2.6) .. (1.6,2.3)
  .. controls (2.3,2.0) and (2.6,2.9) .. (3.1,2.6)
  .. controls (3.7,2.2) and (4.4,1.4) .. (4.9,0.7);
% supporting line: hull chord through the two tangency points, spanning the dent
\draw[red!60!black, thick] (1.0,2.59) -- (5.35,0.48);
\fill[blue!70!black] (1.6,2.3) circle (2pt);
\fill[blue!70!black] (4.9,0.7) circle (2pt);
\fill[red!70!black] (2.56,2.54) circle (2pt);
\draw[-{Latex}] (3.3,4.3) -- (2.64,2.72);
\node[anchor=south, align=center] at (3.4,4.3) {dent: undominated,\\but no line reaches it};
\node[anchor=north, align=center] at (2.7,-0.55) {weighted sum:\\supporting lines, hull only};
\end{scope}
% ---- right panel: chebyshev ----
\begin{scope}[xshift=7.6cm]
\draw[-{Latex}] (0,0) -- (5.4,0) node[below left=1mm and 0mm] {$f_1$};
\draw[-{Latex}] (0,0) -- (0,5.4) node[left] {$f_2$};
\draw[thick, blue!70!black]
  (0.6,4.8) .. controls (0.7,3.4) and (1.0,2.6) .. (1.6,2.3)
  .. controls (2.3,2.0) and (2.6,2.9) .. (3.1,2.6)
  .. controls (3.7,2.2) and (4.4,1.4) .. (4.9,0.7);
% ideal point
\fill (0.6,0.7) circle (1.8pt) node[below right, font=\footnotesize] {$\mathbf{z}^{\star}$};
% corner level sets: vertex slides along the ray from the ideal point
\draw[red!60!black] (0.7,1.8) -- (1.78,1.8) -- (1.78,0.85);
\draw[red!60!black, thick] (1.2,2.54) -- (2.56,2.54) -- (2.56,1.2);
\fill[red!70!black] (2.56,2.54) circle (2pt);
\draw[-{Latex}] (3.6,4.3) -- (2.66,2.7);
\node[anchor=south, align=center] at (3.6,4.3) {growing corners:\\this one reaches the dent};
\node[anchor=north, align=center] at (2.7,-0.55) {Chebyshev~\eqref{eq:chebyshev}:\\corner level sets, whole front};
\end{scope}
\end{tikzpicture}
\caption{Why the obvious scalarization loses solutions (schematic). Left:
minimizing a weighted sum slides a line toward the origin; it can only
stop on the convex hull, so the concave stretch of the front is
unreachable for every choice of weights. Right: the Chebyshev
scalarization's level sets are corners, which can touch any Pareto point,
concave stretches included; sweeping its weights traces the whole front.}
\label{fig:chebyshev}
\end{figure}

Evolutionary methods attack the front directly with a population, using
dominance itself as the fitness signal. NSGA-II (the name abbreviates
non-dominated sorting genetic algorithm), the workhorse, sorts the
population into non-domination ranks (rank 1: the undominated; rank 2:
undominated once rank 1 is removed; and so on) and breaks ties within a
rank by crowding distance, a dispersion measure that prefers points in
sparsely covered stretches of the front, so selection pressure pushes
toward the front while spreading along it \citep{deb2002nsga2}. It needs
generous evaluation budgets, which multi-fidelity variants mitigate, and
it is the default multi-objective sampler in Optuna, so we get a tested
implementation for free when trials are cheap.

\section{Bayesian optimization of the front, one repair at a time}
\label{sec:mobo}

When a trial costs hours, the population's appetite disqualifies it,
and we are back in Chapter~\ref{sec:bayesian}'s situation: few,
expensive, noisy evaluations, and a surrogate that has to make each
one count. The machinery that results is best understood the way EI
was: one natural idea, then a short chain of repairs, each forced by
a failure that is easy to see once it is named.

\subsection{EHVI: expected improvement, paid in area}

The natural idea first. With one objective, EI asked how much a
candidate would improve on the incumbent, averaged over the
surrogate's uncertainty. With several objectives there is no single
incumbent; there is the staircase of Figure~\ref{fig:pareto}, and
progress means adding area to what it dominates. \emph{Expected
hypervolume improvement} (EHVI) therefore scores a candidate by the
expected increase of~\eqref{eq:hv} if we were to evaluate it: the
expected area of fresh territory.

One number makes the definition mechanical. Take the front of
Figure~\ref{fig:pareto} and a candidate configuration whose GP
posterior, say, puts half its mass at the outcome $(3, 1.5)$ and half
at $(3, 2.5)$. The first outcome lands in the dent and carves out the
$1 \times 0.5$ rectangle we computed earlier, an area of $0.5$; the
second is dominated by $(2,2)$ and adds nothing. The candidate's EHVI
is $0.5 \cdot 0.5 + 0.5 \cdot 0 = 0.25$. Every candidate gets such a
score, and the searcher proposes the maximizer, exactly EI's decision
logic with improvement measured in area. In practice the posterior is
a continuous Gaussian rather than two points, so the average is
estimated from posterior samples; nothing about the logic changes.

\subsection{Making it computable in batches: qEHVI}

Turning that definition into a working acquisition hits a geometric
wall. Each posterior sample asks for the area a proposed outcome adds
to a union of overlapping boxes, and we rarely propose one point at a
time: Chapter~\ref{sec:bayesian}'s batch argument applies unchanged,
since idle workers do not care that the objectives are plural. The
joint fresh area of $q$ proposals is the area of a union over a
union, and evaluating it naively by inclusion-exclusion sums over
subsets of the batch, a count that grows exponentially in $q$.

\citet{daulton2020qehvi} made the computation practical at batch
sizes above one. The dominated region is decomposed once per front
into disjoint boxes; each posterior sample's contribution is then a
sum over those pieces, exact rather than approximated, and the whole
estimator is differentiable in the candidate coordinates, so the
batch is optimized by gradient ascent like any other acquisition of
Chapter~\ref{sec:bayesian}.

\subsection{Noise bends the front: qNEHVI}

Now the failure our workloads force. Suppose one Hamiltonian trial,
of genuinely middling quality, draws a lucky seed and reports
$(1.8, 1.8)$. That phantom point dominates the honest compromise at
$(2,2)$, so the observed staircase bends around an outcome the true
objectives never produced. Every genuinely improving candidate is now
scored against territory the phantom already claims, EHVI near zero
across the board, and the study stalls defending a front that does
not exist. This is the winner's curse of
Section~\ref{sec:acquisition} again, with the lucky seed promoted
from incumbent to front member.

qNEHVI applies the same repair noisy EI used there: stop letting raw
observations define the front. The improvement is integrated over the
surrogate posterior at the observed points, so the front itself is
treated as uncertain, and a phantom member holds exactly as much
territory as the posterior believes it deserves
\citep{daulton2021qnehvi}. The same paper removed the remaining cost
wall, making joint batch proposals tractable where the naive
computation grew exponentially in batch size.

\subsection{Vanishing improvements again: qLogNEHVI}

The last repair is inherited. Late in a good study, almost no
candidate adds much area, and expected improvements shrink to the
magnitudes where Section~\ref{sec:acquisition} showed plain EI
underflowing to an exact floating-point zero; the underflow argument
applies verbatim to hypervolume improvements. The fix is the same
log-space computation, and the numerically stabilized variant
qLogNEHVI is the recommended default in the BoTorch ecosystem
\citep{ament2023logei}. The lineage that looked like alphabet soup at
the chapter's door is thus three cumulative repairs to one idea, and
the name now reads as a recipe: the q is the batch, the N is the noise
treatment, the Log is the underflow repair, and EHVI is the idea they
all serve.

\subsection{Two extensions on file}

Two further members solve problems we do not have yet, and one
paragraph each is enough to know when to reach for them. For
high-dimensional spaces, MORBO transplants the trust-region idea of
Figure~\ref{fig:turbo}: several local regions each pursue
hypervolume improvement with local models, coordinating through the
shared front \citep{daulton2022morbo}.

And when the objectives differ in evaluation cost, a profiler's
latency reading next to a full training run's accuracy, the
hypervolume knowledge gradient extends the machinery to decoupled
evaluations: it decides not only where to evaluate next but which
objective to measure there \citep{daulton2023hvkg}.

Our product needs qLogNEHVI as the premium
two-to-four-objective searcher and Chebyshev scalarization as the
cheap many-objective fallback; both ride on the same GP
infrastructure as Chapter~\ref{sec:bayesian}.

\section{Multi-objective meets multi-fidelity}

The savings of Chapter~\ref{sec:multifidelity} apply here too, with one
twist: promotion decisions need a ranking, and ``better'' is now the
partial order we started the chapter with. Promote-the-top-third is not
defined when the top third is not defined.

MO-ASHA resolves this with the layering NSGA-II already taught us.
Rung members are sorted into non-domination ranks, and rank 1 promotes
before rank 2. Ties within a front are broken by a dispersion rule in
objective space, either NSGA-II's crowding distance or an epsilon-net
that repeatedly picks the candidate farthest from those already
chosen; the top fraction then promotes as usual. In its authors'
experiments, promotion rules aware of the whole front beat scalarized
promotion in wall-clock terms \citep{schmucker2021moasha}, which is
the geometry of Figure~\ref{fig:chebyshev} surfacing inside the
scheduler: a scalarized promotion inherits the scalarization's
blind spots, rung by rung.

The repair travels. MO-DEHB applies the same
non-dominated-sorting-plus-crowding selection to DEHB's evolutionary
rungs \citep{awad2023modehb}, and the population family of
Chapter~\ref{sec:population} has its member too: MO-PBT swaps PBT's
scalar exploit ranking for non-dominated selection, evaluated on
trade-offs such as precision against recall and accuracy against
fairness \citep{dushatskiy2023mopbt}.

These are exactly the schedulers we will pair with cheap objectives
such as short-horizon energy drift in Hamiltonian networks. The drift
at rung 1 is cheap to measure and, per Figure~\ref{fig:curves}'s
warning, worth a pilot check that its early rankings mean anything
before we let it cull.

\section{Constraints and preferences}

Two pragmatic variants round out the toolkit. Hard constraints (memory
ceilings, latency budgets, a minimum acceptable return) are handled by
weighting acquisitions with the modeled probability of feasibility. A
candidate whose acquisition score is $0.8$ but whose modeled chance of
fitting under the memory ceiling is $0.3$ counts as
$0.8 \times 0.3 = 0.24$: still on the table, discounted by exactly its
chance of being usable. Noisy-observation treatments of the same idea
exist \citep{letham2019noisyei}, and trust-region versions scale it to
many constraints at once \citep{eriksson2021scbo}.

Preferences short of hard
constraints (``we care about drawdown roughly twice as much as
turnover'') can be encoded by restricting the reference point or the
weight distribution of the scalarization; this is cheap, transparent,
and usually all a practitioner wants. We recommend exposing both, and
resisting anything more elaborate until a real use case demands it.

\section{The evidence, and the check it still lacks}
\label{sec:mo-evidence}

Weighed the way Section~\ref{sec:bo-evidence} taught us to weigh, this
chapter's support has an unusual profile. A larger share than usual is
mathematics that can be checked at a desk; the experiments that exist
are the method authors' own; and the tier that anchored the last two
chapters, an independent budget-matched comparison by outsiders, has no
counterpart here that our literature sweeps could find. Each of those
three facts shapes a different part of the decision box.

Start with what the mathematics settles on its own. The weighted-sum
blind spot is a one-page supporting-hyperplane argument, reproduced in
full above; the Chebyshev form's ability to reach every Pareto point is
its geometric mirror \citep{karl2023moo}; hypervolume strictly
increases exactly when the front advances or spreads, the worked
computation under Figure~\ref{fig:pareto}; and the cost of decomposing
the dominated region grows super-polynomially in the number of
objectives \citep{daulton2021qnehvi}, which draws the two-to-four
boundary. These four facts alone decide the exclusion of weighted
sums, the inclusion of Chebyshev as the fallback, the choice of
progress score, and the boundary past which the premium searcher
gives way. No experiment is being asked to carry those verdicts,
which is fortunate, because experiments of the decisive kind are the
scarce resource here.

The experiment behind the premium pick is the qNEHVI paper's own
evaluation, and its design earns a close look because of one choice:
noise is a first-class experimental axis rather than a footnote. The
comparison spans the modern acquisition families (information-based
PESMO, MESMO, and PFES, the diversity-guided DGEMO, scalarizing
TSEMO and MOEA/D-EGO, Thompson-sampling and plug-in variants of the
hypervolume line) on synthetic and real-world problems, with
observation noise scaled as a percentage of each objective's range
and pushed to twenty percent, a regime the authors note prior noisy
work had probed only at one percent \citep{daulton2021qnehvi}.

Every
method degrades as noise rises; qNEHVI degrades least and leads the
noisy benchmarks in both sequential and batch settings, at
competitive wall time thanks to the cached box decompositions. This
is the authors' own arena, so the selection-bias caution applies, but
two things temper it. The winning move is the same
integrate-over-the-posterior treatment whose single-objective version
carried the noise evidence of Chapter~\ref{sec:bayesian}, arriving
here as transferred machinery rather than a fresh conjecture. And the
axis the authors chose to stress, heavy observation noise, is not a
flattering backdrop chosen for the method; it is the axis our
workloads actually live on, seed-noisy RL returns and
effective-sample-size measurements.

The experiment behind the multi-objective early-stopper is the
MO-ASHA paper's, three tasks chosen to be genuinely unlike one
another: architecture search on a tabular benchmark whose exact
Pareto front is knowable, fairness-accuracy trade-offs on a standard
dataset, and Transformer language-model tuning, all run as
thirty-minute wall-clock races with objectives normalized by
empirical quantiles so that no objective's units dominate by
accident \citep{schmucker2021moasha}. The comparison that matters
inside it: promotion rules that see the whole front (non-dominated
sorting with a dispersion tie-break) against three scalarized
promotion rules, each given a hundred random weight vectors. The
front-aware rules win in wall-clock hypervolume, which is the
experimental echo of the geometry above: a scalarized promotion
inherits the scalarization's blindness, rung by rung.

What this family lacks is the third tier: no independent,
budget-matched, seed-split study of multi-objective tuners on live
training runs surfaced in our audit, nothing that does for this
chapter what \citet{eimer2023hyperparameters} does for
Chapters~\ref{sec:multifidelity} and \ref{sec:population}. The
authors' arenas are also not ours; fairness datasets and tabular
architecture-search benchmarks are a long way from Hamiltonian
rollouts and sampler diagnostics.
The consequence is not paralysis, because the geometry and the
transferred noise machinery carry most of the weight; the
consequence is that our own validation suite has to play the role
the literature has not filled yet, and the falsification criterion
below is written accordingly.

\section{The verdicts, argued from four angles}
\label{sec:mo-angles}

The chapter recommends qLogNEHVI as the premium searcher inside the
two-to-four-objective boundary, Chebyshev scalarization as the cheap
universal fallback, NSGA-II by wrapping for bulk cheap trials, and
MO-ASHA as the early-stopper. Four angles, one verdict.

From geometry: the load-bearing exclusions and boundaries follow
from arguments a reader can verify without trusting a single
benchmark: the hull argument disqualifies weighted sums, the corner
level sets qualify Chebyshev, and the decomposition cost draws the
objective-count boundary. Decisions anchored this way are immune to
the selection-bias worries that haunt benchmark-anchored ones, and
the experiments are needed only for the choices the geometry leaves
open.

From economics: committing to weights before searching is assuming
the exchange rate the study exists to discover; returning a front is
handing the decision-maker the actual menu and letting preference be
exercised when the trade-off is visible and the stakes are known.
The Chebyshev-versus-weighted-sum contrast is the familiar one
between corner-shaped and linear indifference curves, and the dent
in Figure~\ref{fig:chebyshev} is precisely where a linear
indifference curve misreports what is attainable. A reader who
would not let a colleague assume a client's utility function before
seeing the budget set should not let a tuner do it either.

From evidence: the experiments that exist, though author-run, stress
the failure modes that matter to us rather than flattering ones,
heavy observation noise for the acquisition and wall-clock
promotion honesty for the scheduler, and their findings are the
experimental shadows of theorems taught above rather than
free-standing claims. Nothing in this family's record contradicts
any result we verified in the single-objective chapters, and the
missing independent tier is named rather than papered over.

From operations: the whole chapter rides rails already bought. The
premium searcher runs on Chapter~\ref{sec:bayesian}'s GP stack and
inherits its LogEI repair; the evolutionary default arrives by
wrapping Optuna, whose maintenance we audited
(Chapter~\ref{sec:libraries}); the early-stopper is a week of work
on top of our ASHA; and the objective-count boundary is enforceable
in software. The marginal cost of honest multi-objective tuning,
given the rest of the product, is weeks, which changes the
include-or-defer calculus in its favor.

And the falsification criterion, stated in advance: the validation
suite adopts the Hamiltonian study, prediction error against energy
drift, as a pinned two-objective task, run at matched budgets across
seeds with final hypervolume and its uncertainty reported. The
premium searcher must beat random-weight Chebyshev and the wrapped
NSGA-II there; MO-ASHA's cheap early signal, short-horizon drift,
must pass the rank-correlation pilot of
Chapter~\ref{sec:multifidelity} before it is allowed to cull; and if
the cheap fallback matches the premium searcher on our tasks, the
premium tier is demoted to an option, the same rule the
single-objective chapters applied to themselves.

\begin{decisions}
\noindent\textbf{Decisions from this chapter.} Pick by the number of
objectives, the noise level, and the price of a trial.
\begin{itemize}[leftmargin=1.3em, itemsep=2pt]
\item \textbf{qLogNEHVI on BoTorch} (include; the premium searcher;
weeks). For two to four objectives with expensive, noisy trials, which
is precisely the Hamiltonian and finance profile. It is the noise-aware,
numerically repaired member of the hypervolume family, and the BoTorch
ecosystem's own recommended default \citep{daulton2021qnehvi,
ament2023logei}. Beyond four objectives its hypervolume arithmetic
becomes the bottleneck, so the default switches to scalarization there;
treat that as a calibrated cutoff on published complexity results
rather than a prohibition, and make it a configurable threshold whose
override is recorded in the study's audit trail.
\item \textbf{Chebyshev scalarization with random weights} (include;
days). The cheap fallback that composes with every single-objective
engine we own and has better coverage than weighted sums: it can reach
concave stretches of the front that weighted sums cannot
(Figure~\ref{fig:chebyshev}). Completeness under random weights
requires the true ideal point and additional assumptions; the
operational approximation using best-seen values gives better coverage
in practice without the theoretical guarantee. This is also the answer
for five or more objectives.
\item \textbf{Plain weighted sums as a product default} (exclude).
They silently discard the concave stretch of the front; if a user
truly has fixed exchange rates, they have a single-objective problem
and should declare it as one.
\item \textbf{NSGA-II by wrapping Optuna} (include; free). For cheap
trials in bulk, the evolutionary route is a tested default we get by
wrapping, not building.
\item \textbf{MO-ASHA with veto gates} (build; about a week over
ASHA). The multi-objective early-stopper, and the veto machinery
doubles as the correctness gate the sampler workload needs
(Chapter~\ref{sec:workloads}).
\item \textbf{MORBO and the hypervolume knowledge gradient} (defer).
Each answers a situation we do not yet have: very high-dimensional
Pareto search, and objectives measurable separately at different
prices. Revisit on demand.
\end{itemize}
\end{decisions}

\section{What would overturn these verdicts}

qLogNEHVI's tier-1 inclusion depends on V09: if it cannot beat
random-weight Chebyshev scalarization and the wrapped NSGA-II on the
Hamiltonian task by a margin that justifies its complexity, it is
demoted to an option. MO-ASHA's veto gates are tied to V10 and V13:
if cheap fidelity signals (short-horizon drift) fail the
rank-correlation pilot, or if a veto-failing configuration ever gets
promoted in validation, the mechanism is disabled and full-fidelity
MO defaults to the scalarized floor. Weighted-sum exclusion is
reopened if the convexity assumption turns out to hold on our tasks,
which V09's results will show. MORBO's deferral is reconsidered when
trust-region BO is proven and the Pareto-aware TR extension becomes
a marginal cost rather than a second research bet.
